% Isolated post-CP16 research; independently audited reduction.
\subsection{The toral first-image lattice at positive ordinary index}

\begin{lemma}[A toral lattice for the unresolved rank problem]
\label{lem:DE-A3-positive-index-image-candidate}
\status{Candidate; complete independent proof audit pending.}
Let $A$ be a mixed abelian surface over an algebraically closed field
of characteristic two. Let $N$ be ample and suppose its ordinary
polarization index is positive. Put $b=h^0(A,N)$ and
$P=\mathbf P(H^0(A,N))$. Then $b$ is even. There is a normalized
toral invariant derivation $e$, acting on $N$, for which
\[
 W=H^0(A,N)=W_0\oplus W_1,\qquad
 e|W_i=i,\qquad\dim W_0=\dim W_1=b/2.
\]
Write $\mathcal B=W^*\otimes\mathcal O_P(-1)$ and
\[
 Z=\mathbf P(W_0)\cup\mathbf P(W_1),\qquad
 \mathcal I=\ker\bigl(\mathcal B
       \xrightarrow{(s,es)}\mathcal O_P^2\bigr).
\]
Outside $Z$, the sheaf $\mathcal I$ is locally free of rank $b-2$.
If $b/2\ge2$, its reflexive extension has
$\det\mathcal I=\mathcal O_P(-b)$ and $H^0(P,\mathcal I)=0$.
The full coherent Frobenius maps into $\mathcal I$ and preserves it.
When $N=M^{3d}$, the expected sharp relative rank $b-2$ is equivalent
to nonvanishing of the determinant of
$F_P^*\mathcal I\to\mathcal I$, a homogeneous form of degree $b$.
\end{lemma}
\begin{proof}
Positive ordinary index means that the actual finite-flat subgroup
$\mu_2$ in the ordinary multiplicative factor of $A[2]$ is contained
in $\ker\lambda_N$. Its restriction of the commutator pairing is
trivial: the associated homomorphism
$\mu_2\to\mu_2^D=\mathbf Z/2$ is zero.
Thus the isotropic descent theorem applies to
$\pi:A\to B=A/\mu_2$; compare
\cite[Corollary~8.11 and proof, p.~117]
{EdixhovenVanDerGeerMoonenAV}.
Equivalently, the restricted theta group is a commutative
multiplicative-type extension of $\mu_2$ by $\mathbf G_m$.
On character groups its extension splits because the quotient
character group is $\mathbf Z$, which is free.
The resulting splitting supplies an actual $\mu_2$-linearization.
Hence $N=\pi^*H$ for an ample $H$ on $B$.

The torsor algebra is
$\pi_*\mathcal O_A=\mathcal O_B\oplus Q$, where $Q$ is the
nontrivial two-torsion line bundle corresponding to the nonzero
character of $\mu_2$. It is nontrivial because
$H^0(A,\mathcal O_A)=k$. Its multiplication identifies $Q^2$ with
$\mathcal O_B$. The infinitesimal generator of this grading is
the unique nonzero invariant toral vector field $e$ with $e^2=e$.
It acts on the pulled-back bundle by the descended connection.
The two section spaces are $H^0(B,H)$ and $H^0(B,H\otimes Q)$.
Ample abelian-variety vanishing and Riemann--Roch make their
dimensions equal to $\chi(H)=b/2$.

Choose an invariant volume and rational two-basis so that
$e=x\partial_x$ and the volume is $(dx/x)\wedge dy$, as in the
mixed Cartier calculation already proved. Use a frame of $N$
pulled back from $B$; it is killed by $e$.
Cartier extracts the even $x$-powers and odd $y$-powers.
Writing a section coefficient as
$f=a^2+xb^2+yc^2+xyd^2$, one obtains
\[
 C(f^2\omega)=0,\qquad C(f\,e(f)\omega)=0.
\]
The second identity also follows by decomposing $f$ into its two
toral weights: the cross term has weight one and is killed by
Cartier, and the square term is killed because $C(\omega)=0$.
Thus, in the root-removed convention,
$M(s)s=M(s)es=0$. The coherent Frobenius matrix is $M(s)^t$,
so its image annihilates $s,es$, as asserted.

For $s=s_0+s_1$, the vectors $s,es$ are independent precisely
when both $s_0,s_1$ are nonzero. This proves the rank and
codimension assertions. The depth lemma proves reflexivity of
the kernel; the determinant and absence of global sections follow
from its displayed presentation. Since the full coherent model
has the same unit part as the relative model (the ambient top
coherent Frobenius is zero), its stable rank is $b-2$ exactly when
the induced Frobenius on this rank-$b-2$ lattice is generically
bijective. Its determinant lies in
$\operatorname{Hom}(\mathcal O_P(-2b),\mathcal O_P(-b))$,
giving the degree-$b$ form. Nonvanishing is the remaining issue,
not an assertion of this lemma.
\end{proof}

\begin{remark}[The exact first-variation test]
\label{rem:DE-A3-positive-index-variation-candidate}
\status{Candidate reduction; no sharpness theorem claimed.}
At a point outside $Z$ where the induced map on $\mathcal I$
has corank exactly two, let $K$ and $Q$ be its two-dimensional
kernel and cokernel. A tangent direction $t$ gives
\[
 D_t:K\longrightarrow Q.
\]
In a local frame the lowest possible term of the determinant
along that direction is a nonzero fixed scalar times
$\det D_t$. Let $P(s)$ be the rectangular inclusion matrix of a
local frame of $\mathcal I$ into $\mathcal B$, and choose a local
left inverse $L(s)$, so that $L(s)P(s)=1$.
The Frobenius matrix is $L(s)M(s)^tP(s)^{[2]}$. Over dual numbers the
derivative of the last factor is zero; the derivative of the
first factor vanishes on $K$. Consequently $D_t$ is obtained
by restricting $M(t)^t$ and projecting to $Q$.
This is a precise sufficient test for a deformation to the
sharp-rank open. The next lemma proves the required first-corank
statement for an etale double pullback of a sharp positive-index
source. Nonvanishing of $\det D_t$ remains unproved in general.
\end{remark}

\begin{lemma}[The restricted first corank at every sharp double pullback]
\label{lem:DE-A3-positive-index-double-corank-candidate}
\status{Candidate; complete independent proof audit pending.}
Let $\pi:A\to B$ be an etale double isogeny of mixed abelian
surfaces. Let $H$ be ample of positive ordinary index, put
$N=\pi^*H$, and use compatible toral lifts on these bundles.
Suppose that a section $s\in H^0(B,H)$ makes Frobenius on
$I_B=\ker(s,es)\subset H^2(B,H^{-1})$ bijective.
Set $s_A=\pi^*s$ and $b=h^0(A,N)$.
Then on $I_A=\ker(s_A,es_A)$ the first Frobenius rank and stable
Frobenius rank are both $b-4$. Its kernel and cokernel have
dimension two. The full coherent first rank need not equal $b-4$.

For this point, the first-variation map $D_t$ of the preceding
remark vanishes whenever $t$ is pulled back from $H^0(B,H)$.
Consequently it factors through the surjective trace map
\[
 H^0(A,N)\longrightarrow H^0(B,H).
\]
Thus the remaining sufficient sharpness test is the nonvanishing
of a quadratic determinant of a linear map
$H^0(B,H)\to\operatorname{Hom}(k^2,k^2)$.
\end{lemma}
\begin{proof}
Write $V_B=H^2(B,H^{-1})$ and $V_A=H^2(A,N^{-1})$.
First, $s$ and $es$ are independent. Indeed, on the dense
complement of the two weight eigenspaces the full Frobenius has
rank at most $b_B-2$, where $b_B=h^0(B,H)$.
The same inequality holds everywhere by vanishing of its
$(b_B-1)$-minors. At a weight eigenpoint, $I_B$ would have
dimension at least $b_B-1$, contradicting the assumed bijectivity
of its restriction. Pullback preserves this independence, so
both section points lie outside the corresponding bad loci $Z$.
The etale double-torsor algebra has the exact sequence
\[
 0\longrightarrow\mathcal O_B
   \longrightarrow\pi_*\mathcal O_A
   \xrightarrow{\operatorname{Tr}}\mathcal O_B
   \longrightarrow0.
\]
Locally it is the Artin--Schreier algebra
$R[z]/(z^2+z-f)$; its trace sends $1$ to zero and $z$ to one.
Tensoring by $H^{-1}$ and taking cohomology gives
\[
 0\longrightarrow V_B\xrightarrow{\pi^*}V_A
   \xrightarrow{\operatorname{Tr}}V_B\longrightarrow0,
\]
because $H^1(B,H^{-1})=0$ and the surface has no third coherent
cohomology. Frobenius with multiplication by the pulled-back
section commutes with both maps: locally this uses
$\operatorname{Tr}(u^2)=\operatorname{Tr}(u)^2$.

Serre duality identifies trace with the dual of pullback on
sections. Since $es_A=\pi^*(es)$, it follows that
$I_A=\operatorname{Tr}^{-1}I_B$ and there is an exact sequence
of semilinear Frobenius spaces
\[
 0\longrightarrow V_B\xrightarrow{\pi^*}I_A
   \xrightarrow{\operatorname{Tr}}I_B\longrightarrow0.
\]
All kernel-cokernel arguments may equivalently be made after
linearizing these maps from the Frobenius pullbacks of their
domains. Since $k$ is perfect this preserves the dimensions and
exactness; in the first-variation formula the kernel is understood
in this linearized sense.
By hypothesis $F|I_B$ is bijective. Since the full Frobenius on
$V_B$ maps into $I_B$, its image is exactly $I_B$ and its kernel
$K_B$ has dimension two. Also $I_B\cap K_B=0$.
The kernel-cokernel sequence for this diagram gives
\[
 \ker(F|I_A)=\pi^*K_B,\qquad
 \operatorname{coker}(F|I_A)
       \simeq\operatorname{coker}(F|V_B).
\]
These are both two-dimensional.
If an element of $\ker(F|I_A)$ also lies in the image, write it
as $F(x)=\pi^*k$, where $k\in K_B$.
Taking trace and using bijectivity on $I_B$ gives
$\operatorname{Tr}x=0$, so $x=\pi^*v$.
Then $k=F(v)$ lies in $K_B\cap I_B=0$.
Thus the first image is already Frobenius-bijective.
Its dimension is $\dim I_A-2=b-4$.

It remains to verify the assertion on invariant tangent directions.
For $t=\pi^*u$ and $k\in K_B$, the first-variation class is the
class of $\pi^*F_u(k)$, where $F_u$ denotes the full coherent
Frobenius with multiplication by $u$.
Under the preceding cokernel identification it is the class of
$F_u(k)$ in $V_B/I_B$.
The two functionals defining that quotient are $s,es$.
Let $\beta$ be the root-removed Cartier multiplication pairing.
Symmetry and the toral identity give
\[
 \beta(s,u)=\beta(u,s),\qquad
 \beta(es,u)=\beta(s,eu).
\]
The second identity follows in a descended local frame from
$C(e(f)\omega)=0$, since toral weight-one terms have zero Cartier
image. Because $k\in\ker F_s$, it annihilates every
$\beta(s,v)$; hence both $s$ and $es$ annihilate $F_u(k)$.
The first-variation class is therefore zero.

Finally the torsor trace sequence tensored by $H$, together with
$H^1(B,H)=0$, gives
\[
 0\to H^0(B,H)\xrightarrow{\pi^*}H^0(A,N)
       \xrightarrow{\operatorname{Tr}}H^0(B,H)\to0.
\]
This proves the factorization of the linear first-variation map.
No assertion that its quadratic determinant is nonzero follows
from this factorization alone.
\end{proof}
